\section{Numerical implementation and validation}
\label{sec:numerics}

The production implementation is differentiable and uses automatic differentiation to evaluate the finite-basis Jacobian \(J^A{}_i\). Automatic differentiation is a computational method, not a validation. Every reported derivative is checked against symmetric finite differences,
\begin{equation}
    \frac{\partial y^A}{\partial a_i}
    \simeq
    \frac{
      y^A(a_i+\Delta a_i)
      -
      y^A(a_i-\Delta a_i)}
      {2\Delta a_i},
    \label{eq:fd}
\end{equation}
over a range of step sizes.

Before any response-geometry result is used, we require:
\begin{enumerate}
    \item convergence of \(M\) and \(R\) with integration tolerance;
    \item agreement of \(\Lamtwo\) with an independent implementation or published benchmark;
    \item agreement of \(I\) with an independent implementation or published benchmark;
    \item agreement of automatic and finite-difference derivatives;
    \item stability under EOS basis refinement;
    \item repetition with more than one EOS covariance metric;
    \item explicit tests of observable-coordinate dependence.
\end{enumerate}

The coordinate test is important because a numerical measure of thinness can change under nonlinear reparameterizations of observables. We therefore distinguish the physical statement that EOS perturbations approximately preserve a relation from the numerical value of a particular norm assigned to that perturbation.

\todo{Stage 1: insert benchmark table and convergence figure.}
